% Fragmento público generado desde propietarios íntegros.
% Residencia: 52.1.
% Fuente: ../incoming/radion_monografia_20260827/manuscrito/sections/01_resumen_y_guia.tex:1-128
\paragraph{Synthèse mathématique de l'objet et de sa chaîne causale}

Cette monographie reconstruit le radian à partir de la chaîne générative
\[
\APP\longrightarrow\TRIT\longrightarrow\TPK
\longrightarrow X_{\mathrm{enr}}
\longrightarrow\text{structure discrète du continu}
\longrightarrow\MD,
\]
sans prendre comme entrée ni la valeur métrologique de \(180/\pi\), ni le petit
résidu angulaire que l'on souhaite expliquer. Le résultat est une unité enrichie,
le \emph{radion}, défini comme l'unité orientée qui parcourt un radian de
phase et balaie exactement une action \(\hret\) dans l'ellipse positive associée
au même état.

La clôture directrice est
\BigRadionEquation
avec
\[
S_E=2\pih\left(\frac5{36}+\frac{25}{9}\alphah\right),
\qquad
\Delta_4=
\frac{\pih-e_{\HMT}\log\pih}{270}
\left(1+\frac{\pih}{729}\right),
\]
\[
\rho_{\mathrm{tw}}=\frac{2\cdot90}{729}=\frac{20}{81},
\qquad
\epsone=\frac{20}{81}\Delta_4-(S_E-1).
\]
La division APP prouve \(1<S_E<2\), de sorte que le quotient est exactement
un et le reste est \(D_E=S_E-1\). L'égalité
\[
1=S_E-\frac{20}{81}\Delta_4+\epsone
\]
est alors une identité interne. La carte angulaire ultérieure produit
\[
\epsR^\circ
=5.13830167585752820716851646226459474899085744\ldots
\times10^{-8}\,{}^\circ.
\]

Le texte démontre en outre que l'ellipse d'action
\[
Q=a_S\cos\theta,\qquad P=b_S\sin\theta,\qquad
a_Sb_S=2\hret
\]
satisfait
\[
\operatorname{Area}(0,\theta)=\hret\theta.
\]
Par conséquent, un radion balaie \(\hret\) et le tour \(2\pi\) enferme
\(h_{\mathrm{ret}}=2\pi\hret\). La même ellipse possède un bord symplectique
fini et un intérieur d'aire infinie lorsqu'elle est munie de la métrique
de Hilbert--Beltrami--Klein. Telle est la forme exacte de la thèse physique qui guide
l'ouvrage : la surface visible clôt l'action ; la profondeur intérieure
conserve un développement non épuisable par la projection plane.

L'exposition maintient typées quatre réalisations du même état :

\begin{enumerate}[label=\textbf{\arabic*.}]
\item le cercle \(S^1\), qui publie la phase ;
\item l'ellipse positive, qui publie l'action et l'anisotropie ;
\item l'intérieur de Klein, qui mesure la profondeur au moyen du birapport ;
\item l'hélice nonadique, qui conserve retour, mémoire et contraction.
\end{enumerate}

Ce ne sont pas quatre géométries superposées sans critère. Ce sont quatre lecteurs ayant
des domaines, des codomaines et des données conservées différents, appliqués à un unique état
enrichi. Le secteur \(90/120\), la dodécaphase, la retenue du radion, la queue
de Lambert, le noyau de Barbero et les tours globales sont également présentés
comme des objets apparentés mais non interchangeables.

\begin{thesisbox}
Le radion est l'unité orientée qui balaie \(\hret\) dans l'ellipse d'action.
Le cercle publie sa phase, Klein mesure sa profondeur, l'hélice conserve sa
mémoire et la retenue coinductive raccorde exactement ses sections directe et
torsionnelle.
\end{thesisbox}

\paragraph{Ordre d'exposition et dépendances de lecture}

L'ordre n'est pas ornemental. Les Parties I et II produisent l'objet et sa clôture
à partir d'APP--TRIT--TPK. Ce n'est qu'ensuite que sont introduites les réalisations dans le langage
de Hilbert, LC, Maxwell, Minkowski, Hodge, Lorentz, de la modularité et des orbites. Chaque
chapitre distingue trois niveaux :

\begin{description}[style=nextline,leftmargin=4.2cm,labelwidth=3.9cm]
\item[Résultat interne.] Égalité ou construction démontrée dans le domaine
HMT avec ses primitives déclarées.
\item[Formalisation réunie.] Nouvelle application articulant des résultats déjà existants
sans utiliser la cible comme sélecteur.
\item[Interface physique.] Traduction dimensionnelle ultérieure reconnaissant une
structure conventionnelle sans en faire l'origine de la sélection.
\end{description}

Les encadrés de \emph{statut et falsificateur} indiquent quelle observation concrète
invaliderait chaque promotion forte. Cette discipline n'affaiblit pas la thèse ; elle la
rend auditable. L'objectif de la narration étendue est qu'aucun mot
comme \emph{cercle}, \emph{ellipse}, \emph{torsion}, \emph{charge} ou
\emph{mémoire} ne masque un changement de type.

\paragraph{Résultats directeurs et portée probatoire}

\begin{enumerate}[label=\textbf{C\arabic*.}]
\item Clôture indépendante de la valeur finale du radion par deux publications du même état et
soustraction avec retenue APP.
\item Dérivation du coefficient \(20/81\) comme deux feuillets du secteur actif
\(90\), normalisés par la fibre \(3^6=729\).
\item Égalisateur typé entre la couture \(\Re s=1/2\), la face APP de cent
marques et la phase dodécaphasique \(\theta_2=50^\circ\).
\item Théorème d'aire du radion : \(\operatorname{Area}(\theta)=\hret\theta\).
\item Étude focale exacte de l'ellipse et de ses invariants euclidiens,
hyperboliques, symplectiques et modulaires.
\item Théorème bord--intérieur : action totale finie \(h\) et intérieur de Klein
d'aire hyperbolique infinie.
\item Réalisation exacte par covariance minimale et oscillateur LC, suivie de
la publication électromagnétique \(90/120\).
\item Atlas de non-identifications pour tous les objets \(90/120\), y compris
la séparation entre \(\epsR\), queue spectrale et noyau de Barbero \(12:-1\).
\item Caractérisation structurelle de \(\pi\) dans la carte du radion, distincte
du générateur coinductif primaire.
\item Synthèse ontologique : dimension comme rang minimal de mémoires qu'une
projection de phase ne peut plus conserver.
\end{enumerate}
% Fuente: ../incoming/radion_monografia_20260827/manuscrito/sections/01b_mapa_resultados.tex:1-133
\paragraph{Correspondance entre résultats, types et résidences causales}

La correspondance suivante ordonne les affirmations directrices et leurs démonstrations,
en maintenant visible la hiérarchie qui empêche que
la richesse physique efface la clôture mathématique qui la soutient.

% HMT_RESULT_BEGIN:RDN-01-GENEALOGIA:theorem
\begin{exactbox}[title={Généalogie APP--TRIT--TPK du radion}]
Le radion naît de la composition effective
\(\APP\to\TRIT\to\TPK\to X_{\mathrm{enr}}\to\MD\) ; il ne s'obtient pas
en redéfinissant rétrospectivement le radian conventionnel.
\end{exactbox}
% HMT_RESULT_END:RDN-01-GENEALOGIA:theorem

% HMT_RESULT_BEGIN:RDN-02-CIERRE:theorem
\begin{exactbox}[title={Clôture exacte et indépendance prospective}]
La retenue APP produit intérieurement
\[
1=S_E-\frac{20}{81}\Delta_4+\varepsilon_R^{(1)},
\qquad
\varepsilon_R^{(1)}=\frac{20}{81}\Delta_4-(S_E-1),
\]
et la carte angulaire transporte cette même identité vers
\(180^\circ/\pi_{\!\HMT}\) sans introduire le résidu cible.
\end{exactbox}
% HMT_RESULT_END:RDN-02-CIERRE:theorem

% HMT_RESULT_BEGIN:RDN-03-BISAGRA-50:theorem
\begin{exactbox}[title={Coordonnée critique de \(50^\circ\)}]
La coordonnée critique \(\Re s=1/2\), publiée sur la face APP de cent
marques et sur la roue dodécaphasique, sélectionne la phase
\(\theta_2=50^\circ\). C'est une égalité de lecteurs typés du même état.
\end{exactbox}
% HMT_RESULT_END:RDN-03-BISAGRA-50:theorem

% HMT_RESULT_BEGIN:RDN-04-COEFICIENTE-2081:theorem
\begin{exactbox}[title={Densité bidirectionnelle nonadique \(20/81\)}]
Le coefficient torsionnel n'est pas ajusté sur le réseau \(1/81\) : il est construit comme
\(\rho_{\mathrm{tw}}=2N_6/3^6=2\cdot90/729=20/81\), en conservant feuillet,
orientation et normalisation de la fibre nonadique.
\end{exactbox}
% HMT_RESULT_END:RDN-04-COEFICIENTE-2081:theorem

% HMT_RESULT_BEGIN:RDN-05-AREA-ACCION:theorem
\begin{exactbox}[title={Loi surfacique du radion}]
Pour l'ellipse positive
\(Q=a_S\cos\theta\), \(P=b_S\sin\theta\),
\(a_Sb_S=2\hbar_{\mathrm{ret}}\), l'aire orientée balayée est
\(\mathcal A(\theta)=\hbar_{\mathrm{ret}}\theta\). Un radion balaie exactement
\(\hbar_{\mathrm{ret}}\), et le tour \(2\pi\) enferme
\(h_{\mathrm{ret}}=2\pi\hbar_{\mathrm{ret}}\).
\end{exactbox}
% HMT_RESULT_END:RDN-05-AREA-ACCION:theorem

% HMT_RESULT_BEGIN:RDN-06-INVARIANTES-FOCALES:theorem
\begin{exactbox}[title={Invariants focaux de l'opérateur positif}]
Le couple publié \(A\pm C^\ast\) fixe les demi-axes, l'excentricité, la
séparation focale et le paramètre de compression symplectique. En particulier,
\((d_1/d_0)^2=C^\ast/2\), identité adimensionnelle qui sépare la forme de l'échelle.
\end{exactbox}
% HMT_RESULT_END:RDN-06-INVARIANTES-FOCALES:theorem

% HMT_RESULT_BEGIN:RDN-07-BORDE-INTERIOR:theorem
\begin{exactbox}[title={Finitude symplectique du bord et divergence de l'aire hyperbolique intérieure}]
Le bord symplectique de l'ellipse enferme l'action finie
\(h_{\mathrm{ret}}\), tandis que l'aire de Hilbert--Beltrami--Klein de l'intérieur
diverge lorsqu'elle atteint la frontière. Phase visible et profondeur ne sont pas la même
mesure.
\end{exactbox}
% HMT_RESULT_END:RDN-07-BORDE-INTERIOR:theorem

% HMT_RESULT_BEGIN:RDN-08-HELICE-MEMORIA:theorem
\begin{exactbox}[title={Relèvement hélicoïdal nonadique avec retour de phase et avance de mémoire}]
La monodromie nonadique ramène la phase sans réinitialiser l'état. L'hélice
relève le cercle en conservant retenue et mémoire ; la clôture \(2\pi/4\pi\)
distingue retour projectif, inversion d'orientation et retour complet.
\end{exactbox}
% HMT_RESULT_END:RDN-08-HELICE-MEMORIA:theorem

% HMT_RESULT_BEGIN:RDN-09-OSCILADOR-EM:development
\begin{narrativebox}[title={Réalisation oscillatoire et constitutive électromagnétique}]
La même anisotropie se réalise par une covariance minimale et un oscillateur
LC : une coordonnée stocke l'aspect électrique-élastique et sa conjuguée
l'aspect magnétique-rotationnel. La fréquence et l'impédance sont typées
séparément.
\end{narrativebox}
% HMT_RESULT_END:RDN-09-OSCILADOR-EM:development

% HMT_RESULT_BEGIN:RDN-10-TIPOS-90120:theorem
\begin{exactbox}[title={Séparation typée des opérateurs \(90/120\)}]
Extracteur dodécaphasique, retenue du radion, queue de Lambert, fonctionnelle
hexade--octade, noyau de Barbero et semi-groupe global de tours sont des objets
distincts. Le poids \(12:-1\) procède des douze secteurs de la chaîne
fondamentale dodécaphasique et de son unique couture orientée. Le coefficient de
\((1-q)^{-1}\), égal à \(1/24\), et le contre-terme entier minimal vérifient
ultérieurement le couple produit.
\end{exactbox}
% HMT_RESULT_END:RDN-10-TIPOS-90120:theorem

% HMT_RESULT_BEGIN:RDN-11-GEOMETRIAS-CAMPO:development
\begin{narrativebox}[title={Réalisations complexe, lorentzienne et duale de Hodge}]
Les régimes elliptique, parabolique et hyperbolique du TRIT sont publiés comme
rotation, seuil et transformation hyperbolique. L'étoile de Hodge, les lois constitutives et la force
de Lorentz agissent ensuite sur l'état déjà orienté ; elles ne sélectionnent pas ses
constantes.
\end{narrativebox}
% HMT_RESULT_END:RDN-11-GEOMETRIAS-CAMPO:development

% HMT_RESULT_BEGIN:RDN-12-MODULAR-ORBITAL:development
\begin{narrativebox}[title={Caractères modulaire, orbital et arithmétique de frontière}]
L'involution modulaire, les orbites elliptiques, l'avance-retard et l'arithmétique
3--6--9 reconnaissent différentes mémoires du même état. La fonction modulaire et
les vacances sont utilisées comme lecteurs ultérieurs, jamais comme substituts du
générateur APP--TRIT--TPK.
\end{narrativebox}
% HMT_RESULT_END:RDN-12-MODULAR-ORBITAL:development

% HMT_RESULT_BEGIN:RDN-13-SINTESIS-ONTOLOGICA:conclusion
\begin{thesisbox}[title={Synthèse opératoire de phase, d'action, de profondeur et de mémoire}]
Le cercle est la projection de phase ; l'ellipse est la section d'action ; Klein
mesure la profondeur ; l'hélice conserve l'histoire. Le radion est l'unité
orientée qui maintient ces publications couplées sans effacer le résidu qui
les distingue.
\end{thesisbox}
% HMT_RESULT_END:RDN-13-SINTESIS-ONTOLOGICA:conclusion

% HMT_RESULT_BEGIN:RDN-14-CONTROL-ALCANCE:control
\begin{statusbox}[title={Délimitation des domaines et continuité causale}]
Les lois complètes des masses, le raffinement électronique, CKM/CP et le
couloir exceptionnel conservent leurs propres démonstrations directrices. Cette monographie les
cite seulement là où ils partagent un antécédent ou une carte ; elle ne transfère pas au radion leurs
sélecteurs spécifiques et ne confond pas leurs résidences probatoires.
\end{statusbox}
% HMT_RESULT_END:RDN-14-CONTROL-ALCANCE:control
